% ---------------------------------------------------------------------
% Parte 3 · WebSockets · expone Mario · 2:00 (lo edita solo Mario)
% Tres diapositivas: el diagrama, la comparación y un dato con cita.
% ---------------------------------------------------------------------
\begin{frame}{Del HTTP al WebSocket}
  \centering
  \figura{figuras/handshake-websocket.pdf}{\textwidth}{6.0cm}{Diagrama de secuencia propio en
    draw.io (fuente en \texttt{diagramas/}): cliente y servidor, petición de cambio de protocolo,
    respuesta y mensajes en los dos sentidos.}
\end{frame}

\begin{frame}{Qué es WebSocket}
  \begin{columns}[T]
    \begin{column}{0.4\textwidth}
      \begin{tarjeta}{bcmWS}{En una frase}
        \small Es un protocolo estandarizado en 2011. Desde 2013 lo soportan todos los navegadores
        principales, de escritorio y móviles.
      \end{tarjeta}
    \end{column}
    \begin{column}{0.56\textwidth}
      \begin{tarjeta}{bcmTinta}{Qué mirar}
        \small
        \begin{itemize}
          \item El cliente hace un GET normal y pide pasar a WebSocket con la cabecera
            \texttt{Upgrade: websocket}.
          \item Si el servidor acepta, responde con el código 101 y la conexión queda abierta.
          \item Desde ese momento los dos lados mandan mensajes cuando quieren. El servidor ya no espera
            a que le pregunten.
        \end{itemize}
      \end{tarjeta}
    \end{column}
  \end{columns}
  \vfill
  {\scriptsize\color{bcmGris}\cite[secc.~2.4]{Murley2021}\par}
\end{frame}

\begin{frame}{¿Cómo se entera el cliente de un cambio?}
  \begin{columns}[T]
    \begin{column}{0.31\textwidth}
      \begin{tarjeta}[equal height group=alternativas]{bcmGris}{Polling}
        \small
        Pregunta cada cierto tiempo\par\medskip
        Cabeceras HTTP en cada pregunta
      \end{tarjeta}
    \end{column}
    \begin{column}{0.31\textwidth}
      \begin{tarjeta}[equal height group=alternativas]{bcmGris}{Long polling}
        \small
        El servidor espera a tener el dato\par\medskip
        Cabeceras HTTP en cada mensaje
      \end{tarjeta}
    \end{column}
    \begin{column}{0.31\textwidth}
      \begin{tarjeta}[equal height group=alternativas]{bcmWS}{WebSocket}
        \small
        Una conexión que queda abierta\par\medskip
        2 a 14\,B por mensaje
      \end{tarjeta}
    \end{column}
  \end{columns}
  \vfill
  {\scriptsize\color{bcmGris}Mitrović et al.\ (2023) también compararon WebSocket con long polling, en
    este caso en un sistema IoT con ESP32 sobre Wi-Fi \cite[p.~267]{Mitrovic2023}.
    Fuentes: \cite[secc.~2.2 y 2.4]{Murley2021}.\par}
\end{frame}

\begin{frame}{Un dato}
  \begin{columns}[c]
    \begin{column}{0.45\textwidth}
      \centering
      {\color{bcmWS}\fontsize{40}{44}\selectfont\bfseries184,9\,B\par}
    \end{column}
    \begin{column}{0.5\textwidth}
      \small
      \begin{itemize}
        \item Murley y su equipo midieron el tamaño medio de las cabeceras HTTP en los sitios que
          analizaron. Les dio 184,9\,B en las peticiones y 403,1\,B en las respuestas.
        \item En tradingview.com, que actualiza tasas de cambio por WebSocket más o menos cada segundo,
          cada actualización pesa 82\,B en promedio.
        \item Si ese sitio usara long polling, los autores calculan unos 499\,B por actualización.
      \end{itemize}
      {\small\color{bcmGris}\cite[secc.~3.3.4, p.~1197]{Murley2021}\par}
    \end{column}
  \end{columns}
\end{frame}
